\documentclass[11pt,a4paper]{article}
\usepackage[margin=25mm,headheight=15pt]{geometry}
\usepackage[T1]{fontenc}
\usepackage[utf8]{inputenc}
\usepackage{lmodern,amsmath,amssymb,amsthm,mathtools,bm}
\usepackage{xcolor,graphicx,booktabs,tabularx,array,longtable}
\usepackage{microtype,enumitem,fancyhdr}
\usepackage{xurl}
\usepackage{hyperref}
\definecolor{accent}{HTML}{163B55}
\hypersetup{colorlinks=true,linkcolor=accent,citecolor=accent,urlcolor=accent,
 pdftitle={Finite-dimensional detection of reduced almost-periodic Banach cocycles},pdfauthor={Matthew J. Colbrook},
 pdfsubject={AIM 901: mathematical result and proof}}
\urlstyle{same}
\setlength{\parindent}{0pt}
\setlength{\parskip}{0.35em}
\setlength{\emergencystretch}{3em}
\setlist{itemsep=0.25em,topsep=0.3em}
\newtheorem{theorem}{Theorem}[section]
\newtheorem{lemma}[theorem]{Lemma}
\newtheorem{proposition}[theorem]{Proposition}
\newtheorem{corollary}[theorem]{Corollary}
\theoremstyle{remark}
\newtheorem{remark}[theorem]{Remark}
\newtheorem{definition}[theorem]{Definition}
\DeclareMathOperator{\Tr}{Tr}
\DeclareMathOperator{\supp}{supp}
\DeclareMathOperator{\Lip}{Lip}
\DeclareMathOperator{\Rea}{Re}
\newcommand{\R}{\mathbb R}
\newcommand{\C}{\mathbb C}
\newcommand{\E}{\mathbb E}
\newcommand{\N}{\mathbb N}
\newcommand{\ip}[2]{\langle #1,#2\rangle}
\pagestyle{fancy}
\fancyhf{}
\fancyhead[L]{\small\sffamily AIM 901}
\fancyhead[R]{\small\sffamily Matthew J. Colbrook}
\fancyfoot[C]{\thepage}
\renewcommand{\headrulewidth}{0.3pt}
\allowdisplaybreaks[2]
\title{\textbf{Finite-dimensional detection of reduced almost-periodic Banach cocycles}}
\author{Matthew J. Colbrook\\[3pt]
{\small Department of Applied Mathematics and Theoretical Physics}\\
{\small University of Cambridge, United Kingdom}\\
{\small\href{mailto:mjc249@maths.cam.ac.uk}{\texttt{mjc249@maths.cam.ac.uk}}}}
\date{3 November 2026}
\begin{document}
\maketitle
\begin{abstract}
We prove that every nonzero reduced cohomology class of an almost-periodic Banach representation is detected by an unbounded cocycle in a finite-dimensional quotient. Compact-group spectral approximation and separating finite-type quotient maps give the result for arbitrary topological groups.
\end{abstract}
\noindent\textbf{Repository target.}
This manuscript concerns AIM problem 901 at revision
\texttt{5d7e26e8444e}. The precise scope is stated below.
\medskip

\section{Main result and proof}
\label{sec:901}

\paragraph{Answer.} The answer to the stated question is affirmative. In fact,
local compactness of the acting group is unnecessary. The essential distinction
in the proof is between a representation with finitely many irreducible types
and a finite-dimensional representation: the former can have arbitrary Banach
multiplicity spaces.

\begin{theorem}\label{thm:901}
Let $G$ be a topological group, let $V$ be a real or complex Banach space, and
let $\pi:G\to\mathrm{GL}(V)$ be a strongly continuous representation such that
each orbit has compact norm closure. If a continuous $1$-cocycle $b:G\to V$
does not belong to the compact-open closure of
\[
 B^1(G,V)=\{g\mapsto v-\pi(g)v:v\in V\},
\]
then there is a closed invariant subspace $W\subset V$ with
$0<\dim(V/W)<\infty$ such that $\sup_{g\in G}\|b(g)+W\|=\infty$.
\end{theorem}

\subsection{The compact group associated with the representation}

Uniform boundedness gives $M:=\sup_{g\in G}\|\pi(g)\|<\infty$.
Embed $\pi(G)$ into the product of the compact orbit closures, one factor
for each $v\in V$. Every point in its product closure is a linear operator
of norm at most $M$. Consequently its strong operator closure $K$ is compact.
Multiplication is jointly strongly continuous on uniformly bounded sets:
if $A_\alpha\to A$ and $B_\alpha\to B$ strongly with
$\sup_\alpha\|A_\alpha\|<\infty$, then
\[
 \|A_\alpha B_\alpha v-ABv\|
 \le \|A_\alpha\|\,\|B_\alpha v-Bv\|
       +\|(A_\alpha-A)Bv\|\longrightarrow0.
\]
If $\pi(g_\alpha)\to k$, compactness supplies a subnet along which
$\pi(g_\alpha^{-1})\to l\in K$; the displayed estimate implies
$kl=lk=I$. Thus $K$ is a compact group of invertible operators, with
continuous inversion. Write $\rho$ for its strongly continuous action on $V$.
The image of $G$ is dense in $K$.

\subsection{Central spectral approximations}

For now suppose $V$ is complex. For an irreducible unitary representation
$\tau$ of $K$, write $d_\tau=\dim\tau$, $\chi_\tau=\operatorname{tr}\tau$,
and, with normalized Haar measure, define
\[
 E^{\tau}_{ij}=d_\tau\int_K\overline{\tau_{ij}(k)}\rho(k)\,dk,
 \qquad
 P_\tau=\sum_{i=1}^{d_\tau}E^{\tau}_{ii}
       =d_\tau\int_K\overline{\chi_\tau(k)}\rho(k)\,dk .
\]
These formulas mean integration after applying the integrand to a vector.
The resulting operators are bounded. Schur orthogonality and invariance of
Haar measure give the usual matrix-unit relations and
\[
 P_\tau^2=P_\tau,\qquad P_\tau P_\sigma=0\quad(\tau\ne\sigma),
 \qquad
 \rho(k)P_\tau=\sum_{i,j=1}^{d_\tau}\tau_{ij}(k)E^{\tau}_{ij}.
 \tag{901.1}\label{eq:901-matrix}
\]
In particular, all $P_\tau$ commute with $\rho(K)$. On the range of any
finite sum of these projections, the operator span of the representation is
finite dimensional and the representation is continuous in operator norm.
The range itself need not be finite dimensional.

There is a uniformly bounded net $(R_\alpha)$ of finite linear combinations
of the $P_\tau$ such that $R_\alpha v\to v$ for every $v\in V$.
Here is a construction which also addresses nonmetrizable $K$. Compact groups
have a basis of conjugation-invariant neighborhoods of the identity. For
each sufficiently small such neighborhood take a nonnegative continuous
central function supported there, of Haar integral one. Peter--Weyl approximation~\cite[Theorem 1 and Proposition 6]{Tao2011}, followed by averaging under conjugation, approximates this
function in $L^1(K)$ by a central trigonometric polynomial. Choose the
approximation errors tending to zero and the $L^1$ norms at most two.
Integration against $\rho$ gives the asserted net: central trigonometric
polynomials are finite sums of characters, concentration at the identity
gives strong convergence, and $\|R_\alpha\|\le2M$.
It follows, by a finite-net argument, that
\[
 \sup_{v\in C}\|R_\alpha v-v\|\longrightarrow0
 \quad\hbox{for every norm-compact }C\subset V.
 \tag{901.2}\label{eq:901-compact}
\]
Let $P_0=\int_K\rho(k)\,dk$ be the projection onto the invariant vectors.

\subsection{A finite-type cocycle lemma}

\begin{lemma}\label{lem:901-type}
Let $U$ be the range of a finite sum of nontrivial isotypic projections.
Suppose $c:G\to U$ is a continuous cocycle whose image in every
finite-dimensional invariant quotient of $U$ is bounded. Then $c$ is an
exact coboundary.
\end{lemma}

\begin{proof}
First, such quotients separate points of $U$. Indeed, for $\ell\in U^*$,
\[
 D_\ell=\operatorname{span}\{\ell\circ\rho(k)|_U:k\in K\}
\]
is finite dimensional by \eqref{eq:901-matrix}. Its common annihilator
$W_\ell$ is closed and invariant, and $U/W_\ell$ is finite dimensional.
If $u\ne0$, choose $\ell$ with $\ell(u)\ne0$.

The Haar average on $U$ is zero. Since the representation on $U$ is
operator-norm continuous and has finite-dimensional operator span, this
average can be approximated in operator norm by finite convex combinations
of the operators $\pi(g)|_U$. Choose $g_1,\ldots,g_m\in G$ and
$t_j\ge0$, $\sum_jt_j=1$, so that
\[
 A=\sum_{j=1}^m t_j\pi(g_j)|_U,\qquad \|A\|<1,
 \qquad
 v=(I-A)^{-1}\sum_{j=1}^m t_jc(g_j).
 \tag{901.3}\label{eq:901-v}
\]
For any finite-dimensional invariant quotient $q:U\to F$, the bounded
cocycle $qc$ is a coboundary. To see this directly, average an inner
product on $F$ over the compact linear group induced by $K$. The affine
action $z\mapsto\pi_F(g)z+qc(g)$ is then an action by Euclidean isometries
with a bounded orbit. The unique circumcenter of that orbit is fixed.
Thus, for some $w\in F$,
\[
 qc(g)=(I-\pi_F(g))w\qquad(g\in G).
\]
On the quotient, $\|A_F\|\le\|A\|<1$. Applying $q$ to
\eqref{eq:901-v} and using this identity gives $qv=w$. Therefore
\[
 q\bigl(c(g)-(I-\pi(g))v\bigr)=0\qquad(g\in G).
\]
The separating family of finite-dimensional quotients now gives
$c(g)=(I-\pi(g))v$ for every $g$.
\end{proof}

Every quotient in this lemma is available as a quotient of $V$: if $P_U$
is the sum of the corresponding isotypic projections, then
$qP_U:V\to F$ is bounded, onto, and equivariant.

\subsection{Proof of the theorem}

We prove the contrapositive. Suppose that every finite-dimensional invariant
quotient of $V$ sends $b$ to a bounded cocycle. On the trivial isotypic
component, $P_0b$ is a homomorphism into the additive group of $V^K$.
Each scalar functional on this space, composed with $P_0$, is an
equivariant quotient map. A bounded scalar homomorphism is zero, since
$h(g^n)=nh(g)$ for every positive integer $n$. Hence $P_0b=0$.

For a nontrivial type $\tau$, the cocycle $P_\tau b$ satisfies
Lemma~\ref{lem:901-type}, because its finite-dimensional quotient maps
extend to $V$. It is therefore an exact coboundary. Every $R_\alpha b$ is
a finite linear combination of these coboundaries and is itself a
coboundary. If $C\subset G$ is compact, continuity of $b$ makes $b(C)$
norm compact. Equation~\eqref{eq:901-compact} consequently gives
\[
 \sup_{g\in C}\|R_\alpha b(g)-b(g)\|\longrightarrow0.
\]
This proves the complex theorem with the required compact-open closure;
no sequential-closure assumption has been made.

For real $V$, complexify $V$, $\pi$, and $b$. The complexified representation
is still almost periodic: an orbit of $x+iy$ lies in the sum of the compact
orbit closures of $x$ and $iy$. If $b$ were in the compact-open closure of
complex coboundaries, taking real parts of their primitives would put it
in the closure of real coboundaries. Thus a nonzero real reduced class
remains nonzero after complexification. The complex result supplies a
finite-dimensional complex equivariant quotient with unbounded image of
$b$. Restrict that quotient map to $V$ and take its finite-dimensional
real image. Its kernel is the desired $W$. \qed

\paragraph{Scope and source.}
The statement resolves the catalogue's precise finite-quotient question,
associated with Question~4.7 of Cornulier--Tessera~\cite{CT2020}.
The proof uses neither separability nor reflexivity of $V$, nor compact
generation or local compactness of $G$. The uniform approximations have
finite representation type, not finite rank; no approximation property
of the underlying Banach space is assumed. The argument concerns reduced
cohomology and unbounded quotient cocycles, as required in Problem~901.


\begin{corollary}[Consequences for the source questions]
Every locally compact group has property $\mathcal{AP}_{\mathrm{fd}}$
in the terminology of \cite[Definition~4.1]{CT2020}. Consequently
$\mathcal{WAP}_{\mathrm{ap}}$ and $\mathcal{WAP}_{\mathrm{fd}}$ are equivalent
properties of locally compact groups.
\end{corollary}
\begin{proof}
The first assertion is precisely Theorem~\ref{thm:901}. The second
follows from \cite[Proposition~4.3]{CT2020}, which expresses
$\mathcal{WAP}_{\mathrm{fd}}$ as the conjunction of
$\mathcal{WAP}_{\mathrm{ap}}$ and $\mathcal{AP}_{\mathrm{fd}}$.
Thus the same argument answers Questions~4.5--4.7 of that paper;
these consequences are counted as one catalogue entry here.
\end{proof}


\begin{thebibliography}{99}
\small
\setlength{\itemsep}{3pt}

\bibitem{CT2020}
Y.~Cornulier and R.~Tessera,
\emph{On the vanishing of reduced 1-cohomology for Banach representations},
Ann. Inst. Fourier \textbf{70} (2020), 1951--2003.
\href{https://doi.org/10.5802/aif.3363}{doi:10.5802/aif.3363}.
Question 4.7 and Proposition 4.3.

\bibitem{Tao2011}
T.~Tao,
\emph{The Peter--Weyl theorem, and non-abelian Fourier analysis on compact groups},
23 January 2011, Theorem 1 and Proposition 6.
\url{https://terrytao.wordpress.com/2011/01/23/the-peter-weyl-theorem-and-non-abelian-fourier-analysis-on-compact-groups/}.
The group is assumed compact Hausdorff; normalized Haar measure makes
the stated $L^2$ approximation sufficient for $L^1$ approximation.

\end{thebibliography}
\end{document}
